% ============================================================
% ACCESSIBLE PDF SETTINGS
%
% The settings below are important for the final accessibility build,
% but they can make large dissertations compile very slowly.
%
% For routine writing and editing, you may comment out this whole
% \DocumentMetadata block.
%
% Turn it back on when you are:
%   - checking the final structure of the PDF,
%   - testing accessibility,
%   - testing PDF/A-4f or PDF/UA-2 compliance,
%   - preparing the version to submit or share.
%
% I strongly recommend that you address accessibility issues every chapter.
% Do not wait until the thesis is complete to begin accessibility remediation.
% ============================================================

\DocumentMetadata{
  pdfversion=2.0,        % PDF 2.0 output. Keep for final PDF/UA-2 work.
  lang=en-US,            % Sets the document language. Keep for final accessibility work.
  tagging=on,            % SLOW: turns on tagged PDF structure.
  pdfstandard=ua-2,      % Final accessibility target.
  pdfstandard=a-4f,      % Final archival target.
  tagging-setup={
    math/setup={mathml-AF,mathml-SE},
                          % VERY SLOW in large math documents.
                          % Generates MathML tagging for equations.
                          % Important for the final accessible version.
    table/header-rows=1   % Tags the first row of tables as headers.
  }
}

\documentclass[12pt,phd,onehalfspacing]{ncsuthesis}

%% -------------------------------------------------------------------------- %%
%% Core packages for the dissertation template
%% -------------------------------------------------------------------------- %%
%% Keep this list small. Required accessibility and ETD formatting settings are
%% loaded below from ncsu-required.tex.

\usepackage{ncsu-accessibility}
\usepackage{fontspec}
\usepackage{unicode-math}
\usepackage{amsthm}     % Provides the proof environment.
\usepackage{textcomp}
\usepackage{natbib}
%% Include this package to check if text goes over any margin before submission.
%% \usepackage{showframe}
%% Citations should be of the form ``author year'' instead of ``author, year''.
\bibpunct{(}{)}{;}{a}{}{,}

%% -------------------------------------------------------------------------- %%
%% Basic thesis formatting
%% -------------------------------------------------------------------------- %%

\frenchspacing

%% -------------------------------------------------------------------------- %%
%% Student information
%% -------------------------------------------------------------------------- %%
%% Put title, student name, committee members, program, degree year, and PDF
%% metadata fields in student-information.tex.

%%----------------------------------------------------------------------------%%
%% Student and thesis information
%%----------------------------------------------------------------------------%%
%% Edit this file first. The main file, title page, abstract heading, copyright
%% page, and PDF metadata all pull from the values below.

%% Student name.
\newcommand{\NCSUStudentFirstMiddle}{Firstname Middlename}
\newcommand{\NCSUStudentLast}{Lastname}
\newcommand{\NCSUStudentFullName}{\NCSUStudentFirstMiddle\space \NCSUStudentLast}

%% Dissertation or thesis title.
\newcommand{\NCSUThesisTitle}{Title of This Dissertation}

%% Degree program and year.
\newcommand{\NCSUProgramName}{Department of Whichever}
\newcommand{\NCSUDegreeYear}{2026}

%% PDF document properties. These are not printed on the title page, but they
%% help PDF viewers, search tools, and accessibility checkers identify the file.
%% The PDF title and author are taken directly from the title/name fields above.
\newcommand{\NCSUPDFSubject}{NC State University thesis or dissertation}
\newcommand{\NCSUPDFKeywords}{NC State University, thesis, dissertation, accessibility}

%% Committee information.
%% If you have co-chairs, use \cochairI{...} and \cochairII{...} instead of
%% \chair{...}. The committee size includes the chair or co-chairs.
\committeesize{5}
\chair{Chair Person}
\memberI{Member 1}
\memberII{Member 2}
\memberIII{Member 3}
%% If a committee member is an external member, pass in the optional 
%% argument External Member. 
\memberIV[External Member]{Member 4}

%% Send the values above to the NC State thesis class.
\student{\NCSUStudentFirstMiddle}{\NCSUStudentLast}
\program{\NCSUProgramName}
\thesistitle{\NCSUThesisTitle}
\degreeyear{\NCSUDegreeYear}

%% -------------------------------------------------------------------------- %%
%% Required NC State accessible ETD setup
%% -------------------------------------------------------------------------- %%
%% This file contains settings that should remain active for the final
%% accessible NC State thesis/dissertation build.
%%
%% It includes:
%%   - hyperlink and PDF metadata setup,
%%   - PDF bookmark setup,
%%   - microtype,
%%   - NC State Graduate School section/caption formatting,
%%   - final-ETD-safe theorem-like environments.
%%
%% Do not put draft-only design choices in ncsu-required.tex. Those belong in
%% optional-accessible.tex.

%%----------------------------------------------------------------------------%%
%% NC State required accessible ETD setup
%%----------------------------------------------------------------------------%%
%% This file contains settings that should remain active for the final
%% accessible NC State thesis/dissertation build.
%%
%% Do not put draft-only design choices here. Those belong in
%% optional-accessible.tex.
%%----------------------------------------------------------------------------%%

%%----------------------------------------------------------------------------%%
%% Link colors and shared colors
%% The soft colors are defined here because optional-accessible.tex may use
%% them when the optional theorem-card style is enabled.
%%----------------------------------------------------------------------------%%

\usepackage{xcolor}

\definecolor{NCSULinkBlue}{HTML}{005389}
\definecolor{NCSUReynoldsRed}{HTML}{990000}

\definecolor{NCSUSoftBlue}{HTML}{EEF6FB}
\definecolor{NCSUSoftBlueTitle}{HTML}{D9ECF7}
\definecolor{NCSUBlueBorder}{HTML}{2F5D7C}

\definecolor{NCSUSoftGreen}{HTML}{F0F7F2}
\definecolor{NCSUSoftGreenTitle}{HTML}{DCEFE2}
\definecolor{NCSUGreenBorder}{HTML}{3F6F4F}

\definecolor{NCSUBoxGray}{HTML}{F6F6F6}
\definecolor{NCSUBoxGrayTitle}{HTML}{EDEDED}

%%----------------------------------------------------------------------------%%
%% Hyperlinks, PDF bookmarks, and PDF metadata
%%----------------------------------------------------------------------------%%
%% pdfdisplaydoctitle=true tells PDF viewers to display the document title
%% instead of the file name when possible.
%%
%% The visible title page is not enough for accessibility checkers. The PDF also
%% needs a title in the document properties. Do not edit the student name or
%% dissertation title here. Those values are set in student-information.tex.

\usepackage[
  unicode=true,
  bookmarksnumbered=true,
  bookmarksopen=true,
  linktoc=all,
  colorlinks=true,
  allcolors=black,
  pdfdisplaydoctitle=true,
  pdftitle={\NCSUThesisTitle},
  pdfauthor={\NCSUStudentFullName},
  pdfsubject={\NCSUPDFSubject},
  pdfkeywords={\NCSUPDFKeywords},
  pdfcreator={LuaLaTeX with LaTeX PDF management and tagpdf}
]{hyperref}

\usepackage{bookmark}

%%----------------------------------------------------------------------------%%
%% Microtypography
%%----------------------------------------------------------------------------%%
%% Microtype improves spacing and line breaks. This is mainly a readability
%% improvement and should not change the structure of the document.

\usepackage{microtype}

%%----------------------------------------------------------------------------%%
%% NC State Graduate School ETD formatting
%%----------------------------------------------------------------------------%%
%% NC State requires most text to be the same font size as the body text.
%% Section-level heading formatting is defined directly in ncsuthesis.cls so
%% that the template does not depend on the obsolete sectsty package.

\captionsetup{
  font=normalsize,
  labelfont=bf
}

%%----------------------------------------------------------------------------%%
%% Base theorem-like environments
%%----------------------------------------------------------------------------%%
%% The demonstration chapters and many mathematics dissertations use these
%% environments. The base definitions are intentionally plain and suitable
%% for the final ETD:
%%   - no tcolorbox,
%%   - no amsthm theorem-like tagging structure,
%%   - no colored backgrounds.
%%
%% Optional visual changes to these environments belong in optional-accessible.tex.

\newcounter{ncsuStatement}[chapter]
\renewcommand{\thencsuStatement}{\thechapter.\arabic{ncsuStatement}}

\newcommand{\NCSUStatementTitle}[3]{%
  \textbf{#1~#2%
  \if\relax\detokenize{#3}\relax
  \else
    : #3%
  \fi.}%
}

%% Arguments 3 and 4 are reserved for the optional visual style. They are
%% intentionally unused by the plain required style below.
\newcommand{\NCSUStatementBegin}[4]{%
  \refstepcounter{ncsuStatement}%
  \par\addvspace{12pt}%
  \noindent\NCSUStatementTitle{#1}{\thencsuStatement}{#2}\quad
  \begingroup
  \ignorespaces
}

\newcommand{\NCSUStatementEnd}{%
  \par
  \endgroup
  \par\addvspace{12pt}%
}

\DeclareDocumentEnvironment{definition}{O{}}{%
  \NCSUStatementBegin{Definition}{#1}{NCSUSoftGreenTitle}{NCSUGreenBorder}%
}{%
  \NCSUStatementEnd
}

\DeclareDocumentEnvironment{theorem}{O{}}{%
  \NCSUStatementBegin{Theorem}{#1}{NCSUSoftBlueTitle}{NCSUBlueBorder}%
  \itshape
}{%
  \NCSUStatementEnd
}

\DeclareDocumentEnvironment{lemma}{O{}}{%
  \NCSUStatementBegin{Lemma}{#1}{NCSUSoftBlueTitle}{NCSUBlueBorder}%
  \itshape
}{%
  \NCSUStatementEnd
}

\DeclareDocumentEnvironment{proposition}{O{}}{%
  \NCSUStatementBegin{Proposition}{#1}{NCSUSoftBlueTitle}{NCSUBlueBorder}%
  \itshape
}{%
  \NCSUStatementEnd
}

\DeclareDocumentEnvironment{corollary}{O{}}{%
  \NCSUStatementBegin{Corollary}{#1}{NCSUSoftBlueTitle}{NCSUBlueBorder}%
  \itshape
}{%
  \NCSUStatementEnd
}

\DeclareDocumentEnvironment{example}{O{}}{%
  \NCSUStatementBegin{Example}{#1}{NCSUBoxGrayTitle}{black!55}%
}{%
  \NCSUStatementEnd
}

\DeclareDocumentEnvironment{remark}{O{}}{%
  \NCSUStatementBegin{Remark}{#1}{NCSUBoxGrayTitle}{black!45}%
}{%
  \NCSUStatementEnd
}

%% -------------------------------------------------------------------------- %%
%% Student-specific macros and compatibility packages
%% -------------------------------------------------------------------------- %%
%% Put student-specific packages and commands in student-macros.tex.
%% This file should not contain required template settings.
%%
%% This is loaded before optional-accessible.tex so compatibility packages such
%% as accents are in place before optional math fonts are selected.

\IfFileExists{student-macros.tex}{%%----------------------------------------------------------------------------%%
%% Student-specific compatibility macros
%%
%% This file is only for packages and macros needed by this student's chapters.
%%
%% Do not load hyperref, caption, titlesec, trackchanges, tcolorbox, or
%% \begin{document} here. Those belong in the main template files.
%%
%% Theorem-like environments are defined in ncsu-required.tex. Optional visual
%% changes to those environments belong in optional-accessible.tex.
%%----------------------------------------------------------------------------%%

%%----------------------------------------------------------------------------%%
%% Packages needed by the student's content
%%----------------------------------------------------------------------------%%

\usepackage{accents}
\usepackage{comment}
\usepackage{easytable}
\usepackage{multirow}
\usepackage{multicol}

%% This package lets us catch a few literal Unicode characters that may appear
%% in the student's source. It is especially useful for avoiding PDF/A .notdef
%% glyph failures from unsupported symbols.
\usepackage{newunicodechar}

%%----------------------------------------------------------------------------%%
%% TikZ libraries
%%
%% If you do not use TikZ, delete this block.
%% Keep these centralized here so individual figure files do not have to load
%% TikZ libraries themselves.
%%----------------------------------------------------------------------------%%

\usepackage{tikz}
\usetikzlibrary{
  arrows,
  arrows.meta,
  backgrounds,
  calc,
  decorations.pathmorphing,
  decorations.pathreplacing,
  fit,
  math,
  positioning,
  quotes,
  through
}

%% If the original TikZ code used uppercase color names as TikZ keys, these
%% definitions prevent errors such as: I do not know the key '/tikz/Red'.
%%
%% These are RGB-based color aliases, not CMYK definitions.

\definecolor{AccessibleYellow}{HTML}{FFF4B8}
\definecolor{AccessibleCyan}{HTML}{D8F3F5}
\definecolor{AccessibleRed}{HTML}{F7D6D0}
\definecolor{AccessibleGreen}{HTML}{DCEFD8}
\definecolor{AccessibleGray}{HTML}{F2F2F2}
\definecolor{AccessibleLightGreen}{HTML}{EAF6EA}
\definecolor{AccessibleLightBlue}{HTML}{EAF3FA}
\definecolor{AccessibleMediumBlue}{HTML}{BFDDF2}
\definecolor{AccessibleBlue}{HTML}{D9EAF7}
\definecolor{AccessibleLightCyan}{HTML}{D8F3F5}

%%----------------------------------------------------------------------------%%
%% Basic student macros
%%----------------------------------------------------------------------------%%

\providecommand{\uv}[1]{\ensuremath{\mathbf{\hat{#1}}}}
\providecommand{\bo}{\ensuremath{\mathbf{\Omega}}}

\providecommand{\eref}[1]{Eq.~\ref{#1}}
\providecommand{\fref}[1]{Fig.~\ref{#1}}
\providecommand{\tref}[1]{Table~\ref{#1}}
\providecommand{\cref}[1]{Chap.~\ref{#1}}

\providecommand{\todo}[1]{{\color{red}\textbf{TODO: #1}}}

\providecommand{\R}{\mathbb{R}}

%%----------------------------------------------------------------------------%%
%% Other student-specific additions
%%----------------------------------------------------------------------------%%
}{}

%% -------------------------------------------------------------------------- %%
%% Optional visual choices
%% -------------------------------------------------------------------------- %%
%% This file contains optional visual choices only, such as:
%%   - optional chapter heading style,
%%   - optional font choices,
%%   - optional pretty theorem cards.
%%
%% Font choices are loaded here, after student-macros.tex and before luamml.
%% This order allows students to use compatibility packages such as accents
%% while still using optional text-and-math fonts such as Libertinus.
%%
%% You may comment out this file for the most conservative final ETD build.

\IfFileExists{optional-accessible.tex}{%%----------------------------------------------------------------------------%%
%% Optional accessible design choices
%%----------------------------------------------------------------------------%%
%% This file contains optional visual choices only.
%%
%% Required accessibility settings, PDF metadata, hyperlink setup, MathML
%% support, ETD formatting, and student-specific macros belong elsewhere.
%%
%% This file is loaded early in dissertation-main-accessible.tex so that font
%% choices happen immediately after unicode-math. Visual choices are defined
%% here but applied later through \NCSUApplyOptionalVisualChoices.
%%----------------------------------------------------------------------------%%


%% -------------------------------------------------------------------------- %%
%% Optional: thesis font choices
%% -------------------------------------------------------------------------- %%
%% This template is designed for LuaLaTeX, so font changes should use fontspec
%% and unicode-math rather than older LaTeX font packages.
%%
%% Activate at most one font option near the bottom of this section.
%%
%% Leave all font options commented out to use the template's default font setup.

%% -------------------------------------------------------------------------- %%
%% Option 1: Libertinus
%% A polished, readable serif font with a matching math font.
%% -------------------------------------------------------------------------- %%

\newcommand{\NCSUUseLibertinusFonts}{%
  \setmainfont{Libertinus Serif}[
    Ligatures=TeX
  ]%
  \setsansfont{Libertinus Sans}[
    Ligatures=TeX,
    Scale=MatchLowercase
  ]%
  \setmonofont{Libertinus Mono}[
    Scale=MatchLowercase
  ]%
  \setmathfont{Libertinus Math}%
}

%% -------------------------------------------------------------------------- %%
%% Option 2: STIX Two
%% A formal, journal-like font family with strong math support.
%% -------------------------------------------------------------------------- %%

\newcommand{\NCSUUseSTIXTwoFonts}{%
  \setmainfont{STIX Two Text}[
    Ligatures=TeX
  ]%
  \setsansfont{TeX Gyre Heros}[
    Ligatures=TeX,
    Scale=MatchLowercase
  ]%
  \setmonofont{Latin Modern Mono}[
    Scale=MatchLowercase
  ]%
  \setmathfont{STIX Two Math}%
}

%% -------------------------------------------------------------------------- %%
%% Option 3: Latin Modern
%% This is close to the default LaTeX look and is the safest fallback.
%% -------------------------------------------------------------------------- %%

\newcommand{\NCSUUseLatinModernFonts}{%
  \setmainfont{Latin Modern Roman}[
    Ligatures=TeX
  ]%
  \setsansfont{Latin Modern Sans}[
    Ligatures=TeX,
    Scale=MatchLowercase
  ]%
  \setmonofont{Latin Modern Mono}[
    Scale=MatchLowercase
  ]%
  \setmathfont{Latin Modern Math}%
}

%% Activate at most one font option.
%% For the default template font, comment out all three lines.

\NCSUUseLibertinusFonts
% \NCSUUseSTIXTwoFonts
% \NCSUUseLatinModernFonts


%% -------------------------------------------------------------------------- %%
%% Optional: cleaner NC State numbered chapter headings
%% -------------------------------------------------------------------------- %%
%% This changes the visual appearance of numbered chapter headings without
%% changing the actual \chapter command. It avoids decorative chapter-heading
%% packages.
%%
%% It deliberately leaves \chapter* headings unchanged. Therefore, preliminary
%% headings such as ABSTRACT, DEDICATION, BIOGRAPHY, ACKNOWLEDGEMENTS, TABLE OF
%% CONTENTS, LIST OF TABLES, and LIST OF FIGURES retain the official formatting
%% defined in ncsuthesis.cls. Unnumbered headings such as REFERENCES and the
%% APPENDICES divider page also retain the official class formatting.
%%
%% Numbered appendix headings still receive this style because LaTeX implements
%% them as numbered chapters after \appendix.
%%
%% This option is disabled by default so that the official chapter-heading
%% formatting in ncsuthesis.cls remains authoritative. Turn it on only when
%% a deliberately different visual style is desired.

\newif\ifNCSUUseOptionalChapterHeadings

\NCSUUseOptionalChapterHeadingstrue
%\NCSUUseOptionalChapterHeadingsfalse

\makeatletter

\newcommand{\NCSUChapterTitleWithRule}[1]{%
  #1\par\nobreak
  \vspace{0.18in}%
  {\color{black}\rule{0.45\textwidth}{0.8pt}}%
}

\newcommand{\NCSUUseAccessibleChapterHeadings}{%
  \@ifundefined{ParseLaTeXeHeading}{%
    %% Traditional LaTeX chapter driver.
    \renewcommand{\@makechapterhead}[1]{%
      \vspace*{1.0in}%
      {\parindent \z@ \raggedright \normalfont
        \ifnum \c@secnumdepth >\m@ne
          {\normalsize\bfseries\color{black}%
          \MakeUppercase{\@chapapp}\space \thechapter\par}%
          \vspace{0.25in}%
        \fi
        {\Large\bfseries ##1\par}%
        \vspace{0.18in}%
        {\color{black}\rule{0.45\textwidth}{0.8pt}\par}%
        \vspace{0.75in}%
      }%
    }%
  }{%
    %% LaTeX 2026 heading-template instance used by tagged-PDF builds.
    \EditInstance{heading}{ncsu-chapter}{%
      after-penalty-vspace = 1.0in,
      after-vspace         = 0.75in,
      prefix               = \MakeUppercase{\@chapapp},
      heading-decls        = \raggedright\parindent\z@\normalfont,
      prefix-decls         = \normalsize\bfseries\color{black},
      number-decls         = \normalsize\bfseries\color{black},
      title-decls          = \Large\bfseries
    }%
    \EditInstance{headformat}{ncsu-chapter}{%
      heading-indent    = 0pt,
      prefix-number-sep = \WordSpaceAmount{1},
      number-title-sep  = 0.25in,
      title-format      = {\NCSUChapterTitleWithRule{##1}}%
    }%
  }%
}

\makeatother


%% -------------------------------------------------------------------------- %%
%% Optional: pretty theorem cards
%% -------------------------------------------------------------------------- %%
%% The required file defines theorem-like environments in a plain,
%% final-ETD-safe way.
%%
%% This optional block changes the appearance of those environments to a softer
%% math-textbook style. It does not use amsthm or tcolorbox.
%%
%% Turn this on or off with the toggle below.

\newif\ifNCSUPrettyTheorems

% \NCSUPrettyTheoremstrue
\NCSUPrettyTheoremsfalse


%% -------------------------------------------------------------------------- %%
%% Apply optional visual choices
%% -------------------------------------------------------------------------- %%
%% This command is called from dissertation-main-accessible.tex after
%% ncsu-required.tex has been loaded.
%%
%% Do not call this command inside optional-accessible.tex. It must run after
%% ncsu-required.tex because the chapter heading style uses NC State colors and
%% the theorem-card style uses the theorem macros defined there.

\makeatletter

\newcommand{\NCSUApplyOptionalVisualChoices}{%
  \ifNCSUUseOptionalChapterHeadings
    \NCSUUseAccessibleChapterHeadings
  \fi

  \ifNCSUPrettyTheorems
    \@ifundefined{NCSUStatementBegin}{%
      \PackageWarningNoLine{ncsu-template}{%
        Pretty theorem cards were requested, but the base theorem macros were not found.
        Make sure ncsu-required.tex is loaded before optional visual choices are applied.
      }%
    }{%
      \renewcommand{\NCSUStatementBegin}[4]{%
        \refstepcounter{ncsuStatement}%
        \def\NCSUCurrentStatementRuleColor{##4}%
        \par\addvspace{12pt}%
        \noindent
        \begingroup
          \setlength{\fboxsep}{6pt}%
          \colorbox{##3}{%
            \makebox[\dimexpr\linewidth-2\fboxsep\relax][l]{%
              \strut\NCSUStatementTitle{##1}{\thencsuStatement}{##2}%
            }%
          }%
        \endgroup
        \par\nobreak\vspace{6pt}%
        \begingroup
          \leftskip=1em
          \rightskip=1em
          \noindent\ignorespaces
      }%

      \renewcommand{\NCSUStatementEnd}{%
        \par
        \endgroup
        \vspace{4pt}%
        \noindent{\color{\NCSUCurrentStatementRuleColor}\rule{\linewidth}{0.7pt}}%
        \par\addvspace{12pt}%
      }%
    }%
  \fi
}

\makeatother}{}

%% -------------------------------------------------------------------------- %%
%% Apply optional visual choices
%% -------------------------------------------------------------------------- %%
%% The optional file defines visual choices. We apply them here because chapter
%% headings and theorem-card styling depend on commands and colors defined by
%% ncsu-required.tex.

\providecommand{\NCSUApplyOptionalVisualChoices}{}
\NCSUApplyOptionalVisualChoices

%%----------------------------------------------------------------------------%%
%% MathML support
%%----------------------------------------------------------------------------%%
%% This supports the math accessibility settings from \DocumentMetadata.
%%
%% Keep this after optional-accessible.tex so optional font choices, such as
%% Libertinus or STIX Two, are applied before luamml patches the math setup.

\usepackage{luamml}

%% -------------------------------------------------------------------------- %%
%% Begin document
%% -------------------------------------------------------------------------- %%

\begin{document}

%% -------------------------------------------------------------------------- %%
%% Front matter
%% -------------------------------------------------------------------------- %%
%% The front matter contains the pages that come before the main chapters.
%% In this template, that includes the student and thesis information, abstract,
%% dedication, biography, acknowledgements, and the automatically generated
%% lists.
%%
%% The table of contents, list of tables, and list of figures are generated
%% automatically by LaTeX from the chapter titles, section titles, table
%% captions, and figure captions in the thesis.
%%
%% If the table of contents, list of tables, or list of figures looks incomplete,
%% compile the document with LuaLaTeX several times. These lists are built from
%% information saved during earlier compile runs.

\frontmatter

%% ------------------------------ Abstract ---------------------------------- %%
\begin{abstract}

This abstract, as well as the Dedication, Biography, and Acknowledgements below, is being used as a short orientation to this template. Before final submission, replace this entire abstract with the actual dissertation abstract. You can find this section, along with the Dedication, Biography, and Acknowledgements, in the \texttt{front-accessible.tex} file.

These instructions are meant to help you get started gently. You do not need to understand every file in the template before you begin. The main idea is that a thesis or dissertation is too large to keep comfortably in one file, so this project is split into a main file, a front matter file, chapter files, appendix files, a bibliography file, and a few template support files.

The main file is \texttt{dissertation-main-accessible.tex}. This is the file you compile to create your thesis PDF. It controls the document setup, packages, front matter, chapters, bibliography, and appendices. The student name, title, degree program, degree year, committee information, and PDF metadata text are stored in \texttt{student-information.tex}. The front matter is stored in \texttt{front-accessible.tex}. The bibliography entries are stored in \texttt{StudentName-thesis.bib}. Hyperlink settings are stored in \texttt{optional-accessible.tex}.

The main chapters are stored in separate folders, such as \texttt{Chapter-1/Chapter-1.tex}, \texttt{Chapter-2/Chapter-2.tex}, and \texttt{Chapter-3/Chapter-3.tex}. Appendices are stored in separate appendix folders. Keeping this folder structure intact is important. If you rename a file or folder, you must also update the corresponding \texttt{\textbackslash include\{...\}} line in the main file.

This template must be compiled with LuaLaTeX. The ``scholarly content" in the sample chapters is placeholder material only. It is included to test chapter structure, figures, tables, mathematical content, citations, appendices, and accessibility features. Much of the placeholder content was generated with AI.

\end{abstract}

%% ---------------------------- Copyright page ------------------------------ %%
%% Comment the next line if you do not want the copyright page included.
\makecopyrightpage

%% -------------------------------- Title page ------------------------------ %%
\maketitlepage

%% -------------------------------- Dedication ------------------------------ %%
\begin{dedication}
This Dedication is being used for template instructions. Before final submission, replace this text with your own dedication, or remove the dedication if you do not want to include one.

When you begin adapting the template, start with small edits. First, open \texttt{student-information.tex} and update the student name, dissertation title, committee information, program name, degree year, PDF subject, and PDF keywords. Then compile. (Important reminder: This template must be compiled with LuaLaTeX.)

Do not move individual files out of the project folder. The files refer to each other using paths such as \texttt{Chapter-5/Chapter-5.tex}. This means LaTeX expects to find a folder called \texttt{Chapter-5} and a file inside it called \texttt{Chapter-5.tex}. If you move files around, LaTeX may not be able to find them.

The table of contents, list of figures, and list of tables are made automatically by LaTeX. You do not type these pages by hand. Instead, LaTeX builds them from your chapter titles, section titles, figure captions, and table captions. If you change a title or caption and the table of contents does not update right away, compile again. It often takes more than one run for LaTeX to update these lists correctly.

You should usually not edit template support files such as \texttt{ncsuthesis.cls}, \texttt{ncsu-accessibility.sty}, \texttt{changepage.sty}, or \texttt{totcount.sty}. These files control the template, formatting, or helper commands. Most of your work should happen in \texttt{student-information.tex}, \texttt{front-accessible.tex}, the chapter files, the appendix files, and \texttt{StudentName-thesis.bib}.

\end{dedication}

%% -------------------------------- Biography ------------------------------- %%
\begin{biography}

This Biography is being used for accessibility guidance. Before final submission, replace this text with your actual biography.

The goal of this template is to make the final PDF more accessible than a traditional LaTeX thesis or dissertation PDF. The template can help with tagging and structure, but accessibility also depends on the choices you make while writing.

Use real LaTeX structure. Write chapters with \texttt{\textbackslash chapter}, sections with \texttt{\textbackslash section}, and subsections with \texttt{\textbackslash subsection}. Do not simulate headings by manually changing font size or using bold text alone.

Write equations in LaTeX. Do not insert equations as screenshots. Use labels for equations, figures, tables, sections, and chapters that you need to refer to later. Do not manually type equation numbers, figure numbers, or table numbers. You will see examples in the demo Chapter 1 below.

Give every meaningful figure a clear caption and a meaningful text description. The text description should explain the important information in the image for a reader who cannot see it. Do not rely only on color to communicate information. For mathematical figures, describe the mathematical structure, not every decorative detail.

Use real table structures for tables, and make sure tables have clear column headers. Do not use screenshots of tables unless there is no reasonable alternative. Try to keep tables simple, because complicated tables are harder to read visually and harder to make accessible.

If a figure caption contains mathematics, use a math-free optional short caption so the List of Figures does not contain tagged formula content. For example, use

\begin{verbatim}
\caption[Short text caption]{Full caption with \(math.\)}
\end{verbatim}

This allows the full caption under the figure to contain mathematical notation while keeping the List of Figures easier for validation tools and assistive technologies to process.

Automated accessibility checkers are useful, but they do not catch everything. A PDF can pass an automated check and still be hard to use, so manual review still matters.

\end{biography}

%% ----------------------------- Acknowledgements --------------------------- %%
\begin{acknowledgements}

These Acknowledgements are being used for compiling and validation instructions. Before final submission, replace this text with your actual acknowledgements.

\begin{itemize}

\item In TeXShop or another local editor, open \texttt{dissertation-main-accessible.tex} and choose LuaLaTeX before compiling. A full dissertation usually needs more than one compile run. This is normal. LaTeX needs extra runs to build the table of contents, citations, cross-references, list of figures, and list of tables.

\item In Overleaf, set the compiler to LuaLaTeX and set the main document to \texttt{dissertation-main-accessible.tex}. The safest way to begin in Overleaf is to upload the whole project as a zip file. Uploading files one at a time can accidentally flatten the folder structure, which may cause errors when LaTeX tries to find chapter files.

\end{itemize}

For citations, add your bibliography entries to \texttt{bibliography.bib}. If citations appear as question marks, run BibTeX and then run LuaLaTeX again. Without changes to the bibliography, you usually just need to run LuaLaTeX twice to update your document fully. But, with additions to the bibliography, a typical full compile sequence is:

\begin{verbatim}
LuaLaTeX
BibTeX
LuaLaTeX
LuaLaTeX
\end{verbatim}

In TeXShop, this usually means compiling once with LuaLaTeX, then choosing BibTeX and compiling, then choosing LuaLaTeX again and compiling twice more. In Overleaf, you may only need to recompile, but if citations remain unresolved, check that Overleaf is using BibTeX correctly and that the citation keys in your chapters exactly match the keys in \texttt{bibliography.bib}.

If references, citations, figure numbers, table numbers, or the table of contents look wrong, compile again before assuming something is broken. If they still look wrong after several runs, check for misspelled labels or citation keys. (The repeated LuaLaTeX runs are what allow LaTeX to update the table of contents, list of figures, list of tables, citations, and cross-references.)

After compiling the dissertation, check the resulting PDF with \href{https://dev.verapdf-rest.duallab.com/}{VeraPDF}. The browser version is useful for quick checks, but it involves uploading your file to their website which you may not want to do. The \href{https://software.verapdf.org/}{desktop application} runs locally on your machine and is more private.

For this template, the most relevant veraPDF checks are the \texttt{PDF/A-4f validation profile} and the \texttt{PDF/UA-2 + Tagged PDF validation profile}. PDF/A-4f checks archival compliance, while PDF/UA-2 checks accessibility tagging. In veraPDF, you can also leave the Validation Profile as Autodetect and the Output Format as HTML.

The goal is to support more accessible thesis and dissertation PDFs and to aim for WCAG 2.1 AA accessibility principles where they apply. Using this template does not guarantee compliance.

\end{acknowledgements}

\thesistableofcontents
\thesislistoftables
\thesislistoffigures

\chapter*{AUTHORSHIP STATEMENT}
\addcontentsline{toc}{chapter}{Authorship Statement}

List contributions clearly by chapter. Use normal list structure rather than manually bolded paragraphs. Replace this sample authorship statement with the statement required for your dissertation, if applicable.

\begin{description}
  \item[Chapter 1.] Student Name: sole author of Chapter 1.
  \item[Chapter 2.] Student Name: primary writing, coding, analysis, and literature review. Advisor Name: project guidance, editing, and revision.
\end{description}

\vskip 1cm

\noindent\textbf{Use of generative artificial intelligence:} This template was updated by Bevin Maultsby with assistance from ChatGPT. ChatGPT was used to generate content, troubleshoot compilation errors, and edit and refine the instructions. Replace this statement with the statement appropriate for your dissertation, your advisor, your program, and current university guidance.


%% -------------------------------------------------------------------------- %%
%% Main matter
%% -------------------------------------------------------------------------- %%
%% The main matter contains the numbered chapters of the thesis or dissertation.
%% After this point, LaTeX switches to Arabic page numbering and chapter-based
%% numbering.

\mainmatter

%% Include each chapter file below.
%% Each chapter file should usually begin with a \chapter{...} command.
%%
%% You may add, remove, rename, or comment out chapter files as needed.
%% If you rename a chapter folder or file, make sure the path below matches.

\chapter{THE FUNDAMENTAL THEOREM OF LINE INTEGRALS}
\label{chap-one}

\section{Purpose of This Sample Chapter}
\label{sec:sample-purpose}

This chapter is included to show that the template can support ordinary dissertation writing, mathematical exposition, displayed equations, theorem environments, proofs, diagrams, tables, citations, and figures with alt text. A student using this file should replace the mathematical content with their own work, but the structure is intended to model accessible source habits.

In particular, this chapter demonstrates the following practices:
\begin{enumerate}
  \item Use real section and subsection headings instead of visual formatting alone.
  \item Write equations in LaTeX source, not as images.
  \item Give every meaningful figure alternative text in the source file.
  \item Use tables with header rows, not aligned text made with spaces.
  \item Use descriptive captions that explain why the figure or table appears in the chapter.
\end{enumerate}

\section{A First Definition}
\label{sec:line-integral-definitions}

Line integrals give a way to accumulate information along a curve. In vector calculus, one common version measures the component of a vector field in the direction of motion along a parametrized path.

\begin{definition}[Line integral of a vector field]
\label{def:line-integral}
Let \(C\) be a smooth curve parametrized by \(\mathbf r(t)\), where \(a \leq t \leq b\). Let \(\mathbf F\) be a vector field defined on the curve. The line integral of \(\mathbf F\) over \(C\) is
\[
\int_C \mathbf F \cdot d\mathbf r
= \int_a^b \mathbf F(\mathbf r(t)) \cdot \mathbf r'(t)\,dt.
\]
\end{definition}

The notation \(d\mathbf r\) reminds us that the integral follows the direction of the curve. If the parametrization runs from a starting point \(P\) to an ending point \(Q\), then the sign and value of the integral can depend on that direction.

\section{The Fundamental Theorem}
\label{sec:ftli-statement}

The fundamental theorem of line integrals is the vector-calculus analogue of the one-variable fundamental theorem of calculus. It says that if a vector field is the gradient of a scalar function, then the line integral depends only on the endpoint values of that scalar function.

Below, we see two subsections.

\subsection{Statement of the theorem}
\label{sec:ftli-statement-subsection}

\begin{theorem}[Fundamental theorem of line integrals]
\label{thm:ftli}
Let \(f\) be a differentiable scalar-valued function whose gradient is continuous on an open region containing a smooth curve \(C\). Suppose \(C\) is parametrized by \(\mathbf r(t)\), where \(a \leq t \leq b\). Then
\[
\int_C \nabla f \cdot d\mathbf r
= f(\mathbf r(b)) - f(\mathbf r(a)).
\]
\end{theorem}

This result is powerful because it replaces an integral over an entire curve with a calculation at two points. When \(\mathbf F = \nabla f\), the function \(f\) is called a potential function for \(\mathbf F\), and the vector field is called conservative.

\subsection{Proof of the theorem}
\label{sec:ftli-proof-subsection}

\begin{proof}
Since \(C\) is parametrized by \(\mathbf r(t)\), Definition~\ref{def:line-integral} gives
\[
\int_C \nabla f \cdot d\mathbf r
= \int_a^b \nabla f(\mathbf r(t)) \cdot \mathbf r'(t)\,dt.
\]
By the multivariable chain rule,
\[
\frac{d}{dt} f(\mathbf r(t))
= \nabla f(\mathbf r(t)) \cdot \mathbf r'(t).
\]
Therefore,
\[
\int_C \nabla f \cdot d\mathbf r
= \int_a^b \frac{d}{dt} f(\mathbf r(t))\,dt
= f(\mathbf r(b)) - f(\mathbf r(a)),
\]
by the one-variable fundamental theorem of calculus.
\end{proof}

\section{A Diagram of the Idea}
\label{sec:ftli-diagram}

Figure~\ref{fig:ftli-path} shows a curve that begins at \(P\) and ends at \(Q\). The theorem says that when \(\mathbf F = \nabla f\), the integral along the curve is determined by the value of \(f\) at those two endpoints, not by the particular route between them.

\begin{figure}[htbp]
\centering
\begin{tikzpicture}[
  scale=1.05,
  >=Latex,
  alt={A coordinate-plane diagram showing a curved path from point P near the lower left to point Q near the upper right. Several light gray level curves surround the path, and small arrows along the path show its orientation from P to Q.}
]
  \draw[->] (-0.3,0) -- (6.2,0) node[right] {\(x\)};
  \draw[->] (0,-0.3) -- (0,4.5) node[left] {\(y\)};

  \draw[very thick,blue!70!black]
    (0.8,0.65) .. controls (1.6,1.8) and (2.6,0.8) .. (3.4,2.05)
    .. controls (4.0,3.0) and (4.8,2.2) .. (5.3,3.55);

  \fill (0.8,0.65) circle (2.5pt);
  \fill (5.3,3.55) circle (2.5pt);
  \node[below left] at (1.8,0.65) {\(P=\mathbf r(a)\)};
  \node[above right] at (5.3,3.55) {\(Q=\mathbf r(b)\)};

\end{tikzpicture}
\caption[A path in the xy-plane illustrating the major conclusion of the FTLI.]{A path from \(P\) to \(Q.\) For a conservative vector field, the line integral depends only on value of the potential function at the endpoints.}
\label{fig:ftli-path}
\end{figure}

\section{A Short Computation}
\label{sec:ftli-example}

The theorem can be checked in a concrete case. Let
\[
f(x,y)=x^2+y^2,
\]
so that
\[
\nabla f(x,y)=\langle 2x,2y\rangle.
\]
Let \(C\) be parametrized by \(\mathbf r(t)=\langle t,t^2\rangle\), where \(0\leq t\leq 1\). Then \(\mathbf r(0)=\langle 0,0\rangle\) and \(\mathbf r(1)=\langle 1,1\rangle\). By Theorem~\ref{thm:ftli},
\[
\int_C \nabla f \cdot d\mathbf r
= f(1,1)-f(0,0)
= 2.
\]

We can also compute directly from the parametrization:
\[
\begin{aligned}
\int_C \nabla f \cdot d\mathbf r
&= \int_0^1 \nabla f(\mathbf r(t)) \cdot \mathbf r'(t)\,dt \\
&= \int_0^1 \langle 2t,2t^2\rangle \cdot \langle 1,2t\rangle\,dt \\
&= \int_0^1 \left(2t+4t^3\right)\,dt \\
&= \left[t^2+t^4\right]_0^1 \\
&= 2.
\end{aligned}
\]
This agrees with the endpoint calculation.

\section{A Small Table Example}
\label{sec:accessible-table-example}

Tables should be used for data that are genuinely tabular. Table~\ref{tab:ftli-comparison} compares the two approaches used in the example above.

\begin{table}[htbp]
\centering
\caption{Two ways to compute the same conservative line integral.}
\label{tab:ftli-comparison}
\begin{tabular}{p{0.28\textwidth}p{0.52\textwidth}p{0.12\textwidth}}
\toprule
\textbf{Method} & \textbf{Main calculation} & \textbf{Value} \\
\midrule
Endpoint method & Compute \(f(\mathbf r(1))-f(\mathbf r(0))\). & \(2\) \\
Direct integral & Compute \(\int_0^1 \nabla f(\mathbf r(t))\cdot \mathbf r'(t)\,dt\). & \(2\) \\
\bottomrule
\end{tabular}
\end{table}

\section{An Example with Citations}
\label{sec:citations-example}

A dissertation chapter usually includes citations as part of the mathematical or scientific narrative. Parenthetical citations can be written with \verb|\citep|, as in \citep{kiladis2006three}. Narrative citations can be written with \verb|\citet|, as in \citet{tomassini2017interaction}. Students should replace these sample references with sources that are relevant to their own work.

 
\section{A fancy diagram}
\label{sec:chapter-one-fancy-diagram}

Many dissertations in mathematics use diagrams to explain structure. In algebraic topology, algebraic geometry, and homological algebra, commutative diagrams are often used to organize maps between groups, modules, or vector spaces. This chapter gives a sample of how a mathematical dissertation might combine definitions, a theorem, a proof sketch, and an accessible TikZ figure.

The example is the snake lemma. The name comes from the connecting homomorphism that winds through the diagram from a kernel term to a cokernel term.

Commutative diagrams are commonly used to organize algebraic information in topology and homological algebra \citep{sample:diagrammatic-algebra}. The snake lemma is a standard example of a result whose proof is naturally represented by a diagram \citep{sample:snake-lemma}.

\section{Exact Sequences}
\label{sec:exact-sequences}

A sequence of groups and homomorphisms

\begin{equation}
A \xrightarrow{f} B \xrightarrow{g} C
\label{eq:short-sequence}
\end{equation}

is said to be exact at \(B\) if

\begin{equation}
\operatorname{im}(f)=\ker(g).
\label{eq:exact-at-b}
\end{equation}

In words, exactness says that the elements of \(B\) that are sent to zero by \(g\) are exactly the elements that came from \(A\).

A short exact sequence has the form

\begin{equation}
0 \longrightarrow A \xrightarrow{f} B \xrightarrow{g} C \longrightarrow 0.
\label{eq:short-exact-sequence}
\end{equation}

This means that \(f\) is injective, \(g\) is surjective, and the image of \(f\) is the kernel of \(g\).

\section{The Snake Lemma}
\label{sec:snake-lemma}

The snake lemma starts with a commutative diagram whose rows are exact:

\begin{equation}
0 \longrightarrow A' \longrightarrow B' \longrightarrow C'
\end{equation}

and

\begin{equation}
0 \longrightarrow A \longrightarrow B \longrightarrow C.
\end{equation}

There are also vertical maps from the top row to the bottom row. The conclusion of the snake lemma is that these maps produce a new exact sequence involving kernels and cokernels.

\noindent\textbf{Theorem: Snake Lemma.}
Suppose the rows in Figure~\ref{fig:snake-lemma-diagram} are exact and the diagram commutes. Then there is a connecting homomorphism

\begin{equation}
\delta:\ker(\gamma)\longrightarrow \operatorname{coker}(\alpha)
\label{eq:connecting-homomorphism}
\end{equation}

such that the following sequence is exact:

\begin{equation}
\ker(\alpha)
\longrightarrow
\ker(\beta)
\longrightarrow
\ker(\gamma)
\xrightarrow{\delta}
\operatorname{coker}(\alpha)
\longrightarrow
\operatorname{coker}(\beta)
\longrightarrow
\operatorname{coker}(\gamma).
\label{eq:snake-lemma-conclusion}
\end{equation}

\section{The Diagram}
\label{sec:snake-diagram}

Figure~\ref{fig:snake-lemma-diagram} shows the standard snake lemma diagram. The solid arrows show the original commutative diagram. The dashed curved arrow shows the connecting homomorphism, which is the part of the diagram that gives the lemma its name.

\begin{figure}[htbp]
\centering
\begin{tikzpicture}[
  x=1.45cm,
  y=1.15cm,
  >=Latex,
  module/.style={draw, rounded corners=2pt, minimum width=1.05cm, minimum height=0.55cm, align=center},
  zero/.style={minimum width=0.5cm, minimum height=0.4cm, align=center},
  map/.style={->, thick},
  exact/.style={->, thick},
  snake/.style={
    ->,
    very thick,
    dashed,
    decorate,
    decoration={snake, amplitude=1.4mm, segment length=6mm}
  },
  alt={A commutative diagram for the snake lemma. The top exact row is zero to A prime to B prime to C prime. The bottom row is zero to A to B to C. Vertical maps alpha, beta, and gamma go from the top row to the bottom row. Kernel terms appear above the top row and cokernel terms appear below the bottom row. A dashed winding arrow labeled delta connects ker gamma to coker alpha.}
]
  % Main row nodes
  \node[zero]   (z1)  at (-3, 0) {\(0\)};
  \node[module] (Ap)  at (-2, 0) {\(A'\)};
  \node[module] (Bp)  at ( 0, 0) {\(B'\)};
  \node[module] (Cp)  at ( 2, 0) {\(C'\)};

  \node[zero]   (z2)  at (-3,-2) {\(0\)};
  \node[module] (A)   at (-2,-2) {\(A\)};
  \node[module] (B)   at ( 0,-2) {\(B\)};
  \node[module] (C)   at ( 2,-2) {\(C\)};

  % Kernel nodes
  \node[module] (Ka) at (-2, 1.65) {\(\ker(\alpha)\)};
  \node[module] (Kb) at ( 0, 1.65) {\(\ker(\beta)\)};
  \node[module] (Kc) at ( 2, 1.65) {\(\ker(\gamma)\)};

  % Cokernel nodes
  \node[module] (Ca) at (-2,-3.65) {\(\operatorname{coker}(\alpha)\)};
  \node[module] (Cb) at ( 0,-3.65) {\(\operatorname{coker}(\beta)\)};
  \node[module] (Cc) at ( 2,-3.65) {\(\operatorname{coker}(\gamma)\)};

  % Row arrows
  \draw[exact] (z1) -- (Ap);
  \draw[exact] (Ap) -- node[above] {\(f'\)} (Bp);
  \draw[exact] (Bp) -- node[above] {\(g'\)} (Cp);

  \draw[exact] (z2) -- (A);
  \draw[exact] (A) -- node[above] {\(f\)} (B);
  \draw[exact] (B) -- node[above] {\(g\)} (C);

  % Vertical maps
  \draw[map] (Ap) -- node[right] {\(\alpha\)} (A);
  \draw[map] (Bp) -- node[right] {\(\beta\)} (B);
  \draw[map] (Cp) -- node[right] {\(\gamma\)} (C);

  % Kernel maps down to top row
  \draw[map] (Ka) -- (Ap);
  \draw[map] (Kb) -- (Bp);
  \draw[map] (Kc) -- (Cp);

  % Kernel sequence
  \draw[map] (Ka) -- (Kb);
  \draw[map] (Kb) -- (Kc);

  % Cokernel maps from bottom row down
  \draw[map] (A) -- (Ca);
  \draw[map] (B) -- (Cb);
  \draw[map] (C) -- (Cc);

  % Cokernel sequence
  \draw[map] (Ca) -- (Cb);
  \draw[map] (Cb) -- (Cc);

  % Snake connecting homomorphism
  \draw[snake]
    (Kc.south)
    .. controls (2.7,0.8) and (1.4,-1.1) ..
    (0.35,-2.75)
    .. controls (-0.45,-3.65) and (-1.2,-3.65) ..
    (Ca.east)
    node[pos=0.38, right] {\(\delta\)};

  % Labels for rows
  \node[align=center] at (3.45,0) {exact\\top row};
  \node[align=center] at (3.45,-2) {exact\\bottom row};

  % Light guide labels
  \node[align=center] at (-3.2,1.65) {kernels};
  \node[align=center] at (-3.35,-3.65) {cokernels};
\end{tikzpicture}
\caption[A schematic snake lemma diagram]{A schematic snake lemma diagram. The solid arrows show the original commutative diagram with exact rows. The dashed winding arrow represents the connecting homomorphism from \(\ker(\gamma)\) to \(\operatorname{coker}(\alpha)\).}
\label{fig:snake-lemma-diagram}
\end{figure}

\section{Constructing the Connecting Homomorphism}
\label{sec:connecting-homomorphism}

The connecting homomorphism is the key part of the theorem. The construction follows one element through the diagram.

Let \(c'\in \ker(\gamma)\). Since \(c'\in C'\) and the top row is exact, we choose an element \(b'\in B'\) such that

\begin{equation}
g'(b')=c'.
\label{eq:choose-b-prime}
\end{equation}

Now apply the vertical map \(\beta\). Since the diagram commutes,

\begin{equation}
g(\beta(b'))=\gamma(g'(b'))=\gamma(c')=0.
\label{eq:commuting-step}
\end{equation}

Thus \(\beta(b')\in \ker(g)\). Since the bottom row is exact, there is an element \(a\in A\) such that

\begin{equation}
f(a)=\beta(b').
\label{eq:choose-a}
\end{equation}

The image of \(a\) in \(\operatorname{coker}(\alpha)\) is defined to be \(\delta(c')\). In other words,

\begin{equation}
\delta(c')=[a]\in \operatorname{coker}(\alpha).
\label{eq:delta-definition}
\end{equation}

The proof then checks that this definition does not depend on the choices made, that \(\delta\) is a homomorphism, and that the resulting sequence is exact.

\section{Why the Diagram Matters}
\label{sec:why-diagram-matters}

The snake lemma is useful because it turns a large commutative diagram into an exact sequence. In algebraic topology, this kind of result helps compare homology groups that arise from related spaces, pairs of spaces, or chain complexes.

The diagram is not just decorative. It records the logic of the proof:

\begin{enumerate}
  \item Start with an element in \(\ker(\gamma)\).
  \item Lift the element across the top row.
  \item Move down through the vertical map.
  \item Use exactness of the bottom row.
  \item Record the resulting class in \(\operatorname{coker}(\alpha)\).
\end{enumerate}

The winding arrow in the figure summarizes this path through the diagram.

\section{Accessibility Notes for Mathematical Diagrams}
\label{sec:math-diagram-accessibility}

Commutative diagrams can be difficult to interpret with assistive technology. A thesis should not rely on the image alone. The surrounding text should describe the structure of the diagram and explain how the diagram is used in the proof.

For this kind of figure, the alt text should identify the major objects, the direction of the maps, and the role of the special arrow. The caption should explain what the figure represents in the argument. The body text should then walk through the important path in the diagram.

\section{Summary}
\label{sec:chapter-one-summary}

This chapter illustrates how a mathematical dissertation chapter might combine exposition, definitions, theorems, proofs, diagrams, computations, tables, citations, and accessible figures. The mathematical content is an AI-generated sample. The important template lesson is that accessibility starts in the source: headings, equations, captions, alt text, tables, and links should be written carefully before the final PDF is created.

\chapter{Synthesis and Characterization of a Model Organic Compound}
\label{chap-two}

\section{Purpose of This Sample Chapter}
\label{sec:organic-purpose}

This chapter is written as a short example of how a dissertation chapter in organic chemistry might organize a synthesis, reaction optimization, and characterization discussion in an accessible LaTeX document. Here are two sample citations \citep{sample:organic-methods} and \citep{sample:nmr-reporting}.

Organic chemistry chapters often move between narrative explanation, reaction schemes, tabular summaries, and compact characterization data. The accessibility goal is the same as in any other discipline: do not rely on a picture alone. A reaction scheme should have a caption and alt text, a table should use real rows and columns, and important observations from a spectrum should be described in nearby text.

\section{Reaction Design}
\label{sec:reaction-design}

As a model example, consider the reduction of acetophenone to 1-phenylethanol. The reaction is common enough to be recognizable, but simple enough to be used in a template without requiring a specialized chemistry drawing package. Figure~\ref{fig:model-reduction-scheme} gives a simplified reaction scheme.

\begin{figure}[htbp]
\centering
\begin{tikzpicture}[
  >=Latex,
  node distance=2cm,
  alt={A simplified reaction scheme showing acetophenone, labeled compound 1, converted to 1-phenylethanol, labeled compound 2. The arrow is labeled sodium borohydride in methanol, beginning at zero degrees Celsius and warming to room temperature, followed by aqueous workup.}
]
  \node[draw,rounded corners,align=center,minimum width=3.1cm,minimum height=1.2cm] (start) at (0,0)
    {\(\mathrm{PhCOCH_3}\)\\ acetophenone\\ \textbf{1}};
  \node[draw,rounded corners,align=center,minimum width=3.1cm,minimum height=1.2cm] (product) at (7.1,0)
    {\(\mathrm{PhCH(OH)CH_3}\)\\ 1-phenylethanol\\ \textbf{2}};
  \draw[->,very thick] (start) -- node[above,align=center]
    {\(\mathrm{NaBH_4}\), \(\mathrm{MeOH}\)} node[below,align=center]
    {\(0^{\circ}\mathrm{C}\) to room\\temperature;\\ aqueous workup} (product);
\end{tikzpicture}
\caption{A simplified reaction scheme for reducing acetophenone to 1-phenylethanol. A student preparing a real organic chemistry dissertation would usually replace this text-based scheme with a properly drawn chemical structure and an equally descriptive alternative text.}
\label{fig:model-reduction-scheme}
\end{figure}

The reaction is an example of nucleophilic hydride delivery to a ketone. In words, a hydride equivalent adds to the carbonyl carbon, and protonation during workup gives the corresponding secondary alcohol. If the stereochemical outcome matters for the research question, then the chapter should say whether the product is racemic, enantioenriched, or isolated as a single stereoisomer. In this short example, the product is treated as racemic.

\section{General Experimental Approach}
\label{sec:general-experimental-approach}

A typical experimental section gives enough information for a trained chemist to understand what was done and, when appropriate, repeat the experiment. It usually includes the scale, solvent, temperature, atmosphere, purification method, yield, and characterization data. The prose should be compact, but not cryptic.

For example, the following paragraph gives a template-style description of the model reaction. The quantities are included only as examples and should be checked before being used in any actual experiment.

\begin{quote}
A flame-dried round-bottom flask equipped with a magnetic stir bar was charged with acetophenone \textbf{1} (1.00 mmol) and methanol (5 mL). The solution was cooled to \(0^{\circ}\mathrm{C}\), and sodium borohydride (1.20 mmol) was added portionwise. The reaction mixture was allowed to warm to room temperature and was stirred until thin-layer chromatography indicated consumption of the starting material. The reaction was quenched with saturated aqueous ammonium chloride, extracted with ethyl acetate, dried over sodium sulfate, filtered, and concentrated under reduced pressure. Purification by flash column chromatography gave 1-phenylethanol \textbf{2} as a colorless oil.
\end{quote}

This kind of paragraph contains several accessibility concerns. Abbreviations such as TLC should be spelled out on first use. Chemical hazards and specialized procedures should not be hidden in a figure. If color changes, precipitate formation, or chromatographic behavior are important, they should be described in text rather than only shown in a photograph.

\section{Optimization Data}
\label{sec:optimization-data}

Reaction optimization is often summarized in a table. Table~\ref{tab:organic-optimization} gives a small example. The table uses real text, real column headers, and a visible caption. It is not a screenshot of a spreadsheet.

\begin{table}[htbp]
\centering
\caption{Example optimization results for the model reduction of acetophenone. The data are illustrative and are included to show accessible table structure.}
\label{tab:organic-optimization}
\begin{tabular}{p{0.12\textwidth}p{0.18\textwidth}p{0.18\textwidth}p{0.16\textwidth}p{0.18\textwidth}}
\toprule
\textbf{Entry} & \textbf{Hydride source} & \textbf{Solvent} & \textbf{Temperature} & \textbf{Isolated yield} \\
\midrule
1 & \(\mathrm{NaBH_4}\) & Methanol & Room temperature & 82\% \\
2 & \(\mathrm{NaBH_4}\) & Ethanol & Room temperature & 76\% \\
3 & \(\mathrm{NaBH_4}\) & Methanol & \(0^{\circ}\mathrm{C}\) to room temperature & 88\% \\
4 & \(\mathrm{LiAlH_4}\) & Diethyl ether & \(0^{\circ}\mathrm{C}\) & 71\% \\
\bottomrule
\end{tabular}
\end{table}

In the narrative following a table, describe the main pattern rather than forcing readers to infer the conclusion from the table alone. In this example, entry 3 gives the highest isolated yield under the tested conditions. The lower yield in ethanol may reflect slower reduction, while the use of lithium aluminum hydride gives no advantage for this simple substrate.

\section{Characterization Data}
\label{sec:characterization-data}

Characterization data are often compact, but they should still be readable. A student should define unfamiliar abbreviations and explain which data support the structural assignment. For example, proton nuclear magnetic resonance spectroscopy, often abbreviated \({}^{1}\mathrm{H}\) NMR, can distinguish the methyl group and the benzylic proton in 1-phenylethanol. Carbon nuclear magnetic resonance spectroscopy, or \({}^{13}\mathrm{C}\) NMR, gives complementary evidence for the aromatic carbons and the alcohol-bearing carbon.

\begin{table}[htbp]
\centering
\caption{Example characterization summary for 1-phenylethanol. The values are illustrative placeholders for a template and should be replaced with verified experimental values.}
\label{tab:characterization-summary}
\begin{tabular}{p{0.20\textwidth}p{0.66\textwidth}}
\toprule
\textbf{Technique} & \textbf{Example observation} \\
\midrule
\({}^{1}\mathrm{H}\) NMR & Aromatic protons appear near \(\delta 7.2\)--\(7.4\). The benzylic proton appears downfield relative to the methyl group. \\
\({}^{13}\mathrm{C}\) NMR & Aromatic carbon signals and an oxygen-bearing benzylic carbon are observed. \\
Infrared spectroscopy & A broad absorption consistent with an alcohol \(\mathrm{O-H}\) stretch is observed. \\
High-resolution mass spectrometry & The measured mass should be reported with the calculated mass for the proposed molecular formula. \\
\bottomrule
\end{tabular}
\end{table}

A short characterization paragraph might read as follows:

\begin{quote}
Compound \textbf{2} was obtained as a colorless oil. The \({}^{1}\mathrm{H}\) NMR spectrum showed a set of aromatic resonances and signals consistent with a benzylic methine proton and a methyl group. The infrared spectrum contained a broad absorption consistent with an alcohol functional group. These data support assignment of the reduced alcohol product rather than the starting ketone.
\end{quote}

For a real dissertation, the student should replace this paragraph with the full characterization format expected by the discipline, research group, and journal style. The key accessibility point is that the text should state the conclusion supported by the data. A spectrum image may be useful, but it should not be the only place where the structural evidence is communicated.

\section{A Small Stoichiometric Calculation}
\label{sec:stoichiometric-calculation}

Organic chemistry chapters sometimes include short calculations, especially when discussing scale, equivalents, yield, or selectivity. For example, the percent yield is computed from the isolated amount and the theoretical amount:

\begin{equation}
\text{percent yield}=\frac{\text{actual yield}}{\text{theoretical yield}}\cdot 100\%.
\label{eq:percent-yield}
\end{equation}

If the theoretical yield is \(122.17\,\mathrm{mg}\) and the isolated product is \(107.5\,\mathrm{mg}\), then
\begin{align}
\text{percent yield}
  &= \frac{107.5\,\mathrm{mg}}{122.17\,\mathrm{mg}}\cdot 100\% \\
  &= 88.0\%.
\label{eq:percent-yield-example}
\end{align}
This calculation is included as text-based mathematics rather than as an image, so it can be enlarged, searched, copied, and interpreted by assistive technology more reliably than a screenshot.

\section{Accessibility Notes for Chemistry Content}
\label{sec:chemistry-accessibility-notes}

Chemistry documents can be difficult to make accessible because important meaning is often packed into visual structures, spectra, and schemes. The following practices help:
\begin{enumerate}
  \item Give each reaction scheme a caption that states the purpose of the scheme.
  \item Provide alt text that identifies the starting material, product, reagents, and main transformation.
  \item Summarize the important conclusion from each spectrum in text.
  \item Use real tables for optimization and screening data.
  \item Avoid using a screenshot when a table, equation, or short text description would work.
\end{enumerate}

These habits do not replace discipline-specific expectations. They make those expectations easier to meet in a way that is more usable for readers with different access needs.

\chapter{The Central Limit Theorem}
\label{chap-three}

The central limit theorem is one of the main reasons that the normal distribution appears throughout statistics. Roughly speaking, it says that averages of many independent observations tend to look normal, even when the original data do not come from a normal distribution.

This AI-generated chapter gives a short example of how a statistics dissertation might present definitions, a theorem, a proof sketch, a table, and a figure. The goal is model accessible writing in a statistical context.

\section{Sample Means}
\label{sec:sample-means}

Suppose that \(X_1, X_2, \ldots, X_n\) are independent observations from the same population. The sample mean is

\begin{equation}
\bar{X}_n = \frac{1}{n}\sum_{i=1}^{n} X_i.
\label{eq:sample-mean}
\end{equation}

The sample mean is often used to estimate the population mean. If the population mean is \(\mu\), then \(\bar{X}_n\) tends to get closer to \(\mu\) as the sample size increases. This is the content of the law of large numbers.

The central limit theorem gives more information. It describes the shape of the distribution of the sample mean after it has been centered and scaled.

\section{Statement of the Theorem}
\label{sec:clt-statement}

\noindent\textbf{Theorem: Central Limit Theorem.}
Suppose that \(X_1, X_2, \ldots, X_n\) are independent and identically distributed random variables with mean \(\mu\) and finite, positive variance \(\sigma^2\). Then

\begin{equation}
\frac{\bar{X}_n-\mu}{\sigma/\sqrt{n}}
\xrightarrow{d} Z,
\label{eq:clt}
\end{equation}

where \(Z\) has the standard normal distribution.

In words, the standardized sample mean becomes approximately standard normal as \(n\) becomes large. The notation \(\xrightarrow{d}\) means convergence in distribution.

\section{Interpretation}
\label{sec:clt-interpretation}

The expression in Equation~\ref{eq:clt} has three important pieces.

\begin{enumerate}
  \item The numerator \(\bar{X}_n-\mu\) centers the sample mean at zero.
  \item The denominator \(\sigma/\sqrt{n}\) rescales the sample mean by its standard deviation.
  \item The conclusion says that this centered and scaled quantity has an approximately normal distribution for large \(n\).
\end{enumerate}

This result is useful because it allows researchers to use normal approximations in many settings where the original data are not normally distributed.

For example, suppose that the original observations are right-skewed. Individual observations may not look symmetric, but the averages of many such observations may be much closer to normal.

\section{A Proof Sketch}
\label{sec:clt-proof-sketch}

A full proof of the central limit theorem requires more probability theory than we want to include in this sample chapter. However, the main idea can be explained using moment generating functions.

Let

\begin{equation}
Y_i = \frac{X_i-\mu}{\sigma}.
\label{eq:standardized-observation}
\end{equation}

Then each \(Y_i\) has mean \(0\) and variance \(1\). The standardized sample mean can be written as

\begin{align}
\frac{\bar{X}_n-\mu}{\sigma/\sqrt{n}}
&= \frac{\sqrt{n}}{\sigma}\left(\bar{X}_n-\mu\right) \\
&= \frac{\sqrt{n}}{\sigma}
   \left(
     \frac{1}{n}\sum_{i=1}^{n}X_i-\mu
   \right) \\
&= \frac{1}{\sqrt{n}}\sum_{i=1}^{n}
   \frac{X_i-\mu}{\sigma} \\
&= \frac{1}{\sqrt{n}}\sum_{i=1}^{n}Y_i.
\label{eq:standardized-mean}
\end{align}

The proof then studies the distribution of the sum in Equation~\ref{eq:standardized-mean}. Under the hypotheses of the theorem, the moment generating function of this standardized sum approaches the moment generating function of a standard normal random variable. Since moment generating functions determine distributions under appropriate conditions, the standardized sample mean converges in distribution to a standard normal random variable.

\section{A Simulation Example}
\label{sec:clt-simulation}

Table~\ref{tab:clt-simulation} shows a schematic example of what one might observe in a simulation. The original data are assumed to come from a skewed population. As the sample size increases, the distribution of the sample mean becomes more symmetric and more concentrated around the population mean.

\begin{table}[htbp]
\caption{A schematic summary of how sample means behave as the sample size increases.}
\label{tab:clt-simulation}
\centering
\begin{tabular}{lll}
\toprule
Sample size & Shape of sample mean distribution & Variability \\
\midrule
\(n=5\)   & Still noticeably skewed      & Relatively large \\
\(n=30\)  & More symmetric               & Smaller \\
\(n=100\) & Close to normal in many cases & Much smaller \\
\bottomrule
\end{tabular}
\end{table}

The table is not meant to replace a formal theorem. Instead, it gives a practical interpretation of the theorem: larger samples tend to produce sample means whose distributions are more normal and less spread out.

\section{A Diagram of the Normal Approximation}
\label{sec:normal-approximation-diagram}

Figure~\ref{fig:normal-approximation} shows the standard normal curve that appears in the conclusion of the central limit theorem. In an accessible document, the caption identifies the figure, and the surrounding text explains why it matters.

\begin{figure}[htbp]
\centering
\begin{tikzpicture}[
  x=1.25cm,
  y=4cm,
  alt={A bell-shaped standard normal curve centered at zero. The horizontal axis is labeled z, with tick marks at negative two, negative one, zero, one, and two.}
]
  \draw[->] (-3.4,0) -- (3.4,0) node[right] {\(z\)};
  \draw[->] (0,0) -- (0,0.48) node[above] {density};

  \draw[thick, domain=-3.2:3.2, samples=100]
    plot (\x,{0.3989*exp(-0.5*\x*\x)});

  \foreach \x in {-2,-1,0,1,2}
    \draw (\x,0.01) -- (\x,-0.01) node[below] {\(\x\)};

  \node[above] at (0,0.3989) {standard normal curve};
\end{tikzpicture}
\caption{The standard normal curve used in the central limit theorem. After centering and scaling, the sample mean is approximately distributed according to this curve when the sample size is large.}
\label{fig:normal-approximation}
\end{figure}

\section{Why This Matters}
\label{sec:clt-why-it-matters}

The central limit theorem supports many common statistical procedures. For example, confidence intervals and hypothesis tests for means often rely on the fact that standardized sample means are approximately normal.

A typical confidence interval for a population mean has the form

\begin{equation}
\bar{X}_n \pm z_{\alpha/2}\frac{\sigma}{\sqrt{n}},
\label{eq:confidence-interval}
\end{equation}

when \(\sigma\) is known and the sample size is large. Here, \(z_{\alpha/2}\) is chosen from the standard normal distribution.

In practice, \(\sigma\) is often unknown and must be estimated from the data. That leads to related procedures involving the sample standard deviation and the \(t\)-distribution. The central idea is the same: the sample mean is random, but its distribution can often be approximated well enough to make statistical inference possible.

\section{Examples of Image Inclusion}
\label{sec:image-inclusion-examples}

The following figure is intentionally left from the original sample file. It is included here to demonstrate the accessibility helper commands for ordinary image files and cropped image files. 

\accessiblefigurecropped[htbp][\textwidth]
  {4cm 8.5cm 3cm 0cm}
  {Chapter-1/figs/gfgnge_vs_year.eps}
  {A cropped chart from the original template. The student should describe the chart type, axes, trend, and main message when replacing this placeholder.}
  {Sample cropped image included with the original template. Replace this image, alt text, and caption with content from the dissertation.}
  {fig:gfgn}

\section{Accessibility Notes for Statistical Writing}
\label{sec:clt-accessibility-notes}

Statistical writing may include formulas, tables, plots, and notation-heavy explanations. To make this material more accessible, authors should follow these practices.

\begin{enumerate}
  \item Explain the meaning of each major equation in nearby text.
  \item Use real LaTeX equations rather than screenshots of equations.
  \item Use real tables with column headers rather than images of tables.
  \item Describe the purpose and main message of each plot in the surrounding text.
  \item Avoid relying on color alone to communicate statistical conclusions.
\end{enumerate}

These practices help readers who use screen readers, readers who enlarge text, and readers who are new to the notation.
\chapter{Marine, Earth, and Atmospheric Science}
\label{chap-four}

Marine, earth, and atmospheric science often connects observations from the ocean, atmosphere, and land surface. A dissertation in this area may include maps, satellite images, time series, model output, field measurements, and physical equations. This sample chapter shows how that kind of material can be presented in an accessible way.

The example focuses on tropical weather systems. The goal is not to give a complete scientific treatment, but to demonstrate how a student might combine text, figures, tables, and equations in a thesis or dissertation.

\section{Scientific Context}
\label{sec:meas-context}

Tropical cyclones are influenced by several environmental conditions, including sea surface temperature, atmospheric moisture, vertical wind shear, and large-scale circulation patterns. These variables interact across different spatial scales. For example, warm ocean water can provide energy for convection, while strong vertical wind shear can disrupt storm organization.

Figure~\ref{fig:damage} shows an example of a weather-related image from the original template. In a real dissertation, the caption should explain the role of the figure in the argument, and the alt text should describe the visual information needed by a reader who cannot see the image.

\accessiblefigure[htbp][\textwidth]
  {Chapter-1/figs/Damage.png}
  {A map-based hurricane track graphic with labels for Newfoundland, Bermuda, and an African easterly wave, along with inset photographs showing storm damage, flooding, wind, and a damaged boat.}
  {Sample weather-related image included with the original template. In a real dissertation, replace this image, alt text, and caption with content from the dissertation.}
  {fig:damage}

\section{Environmental Variables}
\label{sec:meas-variables}

A typical analysis may compare several environmental variables before, during, and after a storm event. Table~\ref{tab:meas-variables} gives a simple example. The table uses real text and column headers rather than a screenshot.

\begin{table}[htbp]
\caption{Examples of environmental variables used in tropical weather analysis.}
\label{tab:meas-variables}
\centering
\begin{tabular}{lll}
\toprule
Variable & Common interpretation & Example unit \\
\midrule
Sea surface temperature & Ocean energy available to the atmosphere & \(^{\circ}\mathrm{C}\) \\
Vertical wind shear & Change in wind with height & \(\mathrm{m/s}\) \\
Relative humidity & Atmospheric moisture content & percent \\
Surface pressure & Strength of low-pressure circulation & \(\mathrm{hPa}\) \\
Precipitation rate & Rainfall intensity & \(\mathrm{mm/hr}\) \\
\bottomrule
\end{tabular}
\end{table}

Tables like this are useful for readers because they define variables before the variables appear in formulas or figures.

\section{A Basic Physical Relationship}
\label{sec:meas-equation}

One common approximation in atmospheric science is hydrostatic balance. It relates the vertical pressure gradient to the density of air and the acceleration due to gravity:

\begin{equation}
\frac{\partial p}{\partial z} = -\rho g.
\label{eq:hydrostatic-balance}
\end{equation}

In Equation~\ref{eq:hydrostatic-balance}, \(p\) is pressure, \(z\) is height, \(\rho\) is air density, and \(g\) is the acceleration due to gravity. The negative sign indicates that pressure decreases as height increases.

This equation is not a complete model of the atmosphere. Instead, it is a useful approximation when vertical accelerations are small compared with the force of gravity.

\section{A Schematic Cross Section}
\label{sec:meas-cross-section}

Figure~\ref{fig:storm-cross-section} gives a simple schematic of a tropical system over warm ocean water. The figure is drawn with TikZ so that the source remains editable. The alt text describes the main visual features of the diagram.

\begin{figure}[htbp]
\centering
\begin{tikzpicture}[
  x=1cm,
  y=1cm,
  alt={A schematic cross section of a tropical weather system over the ocean. Warm ocean water is shown at the bottom. Curved arrows rise from the ocean surface into a central convective cloud, and upper-level arrows spread outward near the top of the cloud.}
]
  % ocean
  \fill[blue!20] (-5,-1.2) rectangle (5,-0.5);
  \draw[blue!60!black, thick] (-5,-0.5) -- (5,-0.5);
  \node[blue!60!black] at (3.6,-0.85) {warm ocean surface};

  % cloud
  \fill[gray!20] (-1.4,-0.2) .. controls (-2.2,1.1) and (-1.2,2.8) .. (-0.4,3.0)
    .. controls (0.0,3.7) and (1.1,3.4) .. (1.0,2.7)
    .. controls (2.0,2.4) and (2.2,1.2) .. (1.3,0.7)
    .. controls (1.4,0.0) and (0.2,-0.2) .. (-1.4,-0.2);
  \draw[gray!70!black, thick] (-1.4,-0.2) .. controls (-2.2,1.1) and (-1.2,2.8) .. (-0.4,3.0)
    .. controls (0.0,3.7) and (1.1,3.4) .. (1.0,2.7)
    .. controls (2.0,2.4) and (2.2,1.2) .. (1.3,0.7)
    .. controls (1.4,0.0) and (0.2,-0.2) .. (-1.4,-0.2);
  \node at (0,1.4) {deep convection};

  % rising motion
  \draw[->, thick] (-2.8,-0.45) .. controls (-1.8,0.2) and (-1.2,0.8) .. (-0.8,1.6);
  \draw[->, thick] (2.8,-0.45) .. controls (1.8,0.2) and (1.2,0.8) .. (0.8,1.6);
  \draw[->, thick] (0,-0.45) -- (0,2.5);

  % outflow
  \draw[->, thick] (-0.2,3.1) .. controls (-1.2,3.6) and (-2.4,3.7) .. (-3.5,3.4);
  \draw[->, thick] (0.2,3.1) .. controls (1.2,3.6) and (2.4,3.7) .. (3.5,3.4);
  \node at (0,3.95) {upper-level outflow};

  % labels
  \node[align=center] at (-3.7,1.2) {moist air\\converges};
  \node[align=center] at (3.7,1.2) {moist air\\converges};
\end{tikzpicture}
\caption{A schematic cross section of a tropical weather system. Warm ocean water supports rising moist air, deep convection, and upper-level outflow.}
\label{fig:storm-cross-section}
\end{figure}

\section{Example Analysis}
\label{sec:meas-analysis}

Suppose a researcher wants to compare the environment before and during a storm event. One simple approach is to compute changes in variables over time. For example, if \(T_0\) is the sea surface temperature before the event and \(T_1\) is the sea surface temperature during the event, then the temperature change is

\begin{equation}
\Delta T = T_1 - T_0.
\label{eq:temperature-change}
\end{equation}

A negative value of \(\Delta T\) indicates cooling, while a positive value indicates warming. In a real analysis, this calculation might be repeated across many grid points, buoys, or satellite retrievals.

For a spatially averaged value over a region \(R\), one might write

\begin{equation}
\overline{T}
=
\frac{1}{A(R)}
\int_R T(x,y)\,dA,
\label{eq:spatial-average}
\end{equation}

where \(A(R)\) is the area of the region and \(T(x,y)\) is the temperature at location \((x,y)\). The surrounding text should explain the region, the data source, and the units.

\section{Accessibility Notes for Earth Science Figures}
\label{sec:meas-accessibility}

Earth science dissertations often rely on highly visual material. Maps, radar images, satellite images, and model-output plots can be difficult to interpret without careful explanation. Authors should follow these practices.

\begin{enumerate}
  \item Describe the main visual message of each map or plot in nearby text.
  \item Include alt text for every meaningful figure.
  \item Avoid relying on color alone to distinguish categories or intensities.
  \item Use readable labels, legends, and units.
  \item Do not use screenshots of tables or equations.
  \item Explain abbreviations such as SST, RH, MSLP, or AEW the first time they appear.
\end{enumerate}

These practices help readers understand the scientific argument even when they cannot rely on the visual display alone.
\chapter{Computational Methods in Engineering}
\label{chap:computational-methods}

Many dissertations in computer science, engineering, applied mathematics, and the physical sciences include algorithms, source code, numerical experiments, or computational workflows. This sample chapter shows one way to present code in a thesis while keeping the material readable and accessible.

The main example is a finite difference method for the one-dimensional heat equation. The goal is not to provide a complete numerical analysis, but to demonstrate how a student might present a model problem, an algorithm, a short Python script, a short MATLAB script, and a table of computational results.

\section{A Model Problem}
\label{sec:computational-model-problem}

Consider the one-dimensional heat equation

\begin{equation}
\frac{\partial u}{\partial t}
=
\alpha
\frac{\partial^2 u}{\partial x^2},
\label{eq:heat-equation}
\end{equation}

where \(u(x,t)\) is the temperature at position \(x\) and time \(t\), and \(\alpha\) is the thermal diffusivity.

For this demonstration, suppose the spatial domain is \(0\leq x\leq 1\), and suppose the boundary temperatures are fixed:

\begin{equation}
u(0,t)=0
\qquad\text{and}\qquad
u(1,t)=0.
\label{eq:heat-boundary-conditions}
\end{equation}

We also choose the initial temperature distribution

\begin{equation}
u(x,0)=\sin(\pi x).
\label{eq:heat-initial-condition}
\end{equation}

This problem is simple enough to explain in a template, but it still looks like the kind of model problem that appears in engineering and scientific computing.

\section{Finite Difference Approximation}
\label{sec:finite-difference-approximation}

A finite difference method replaces derivatives with difference quotients. If \(x_j=j\Delta x\) and \(t_n=n\Delta t\), then we write \(u_j^n\) for an approximation to \(u(x_j,t_n)\).

A common explicit approximation is

\begin{equation}
u_j^{n+1}
=
u_j^n
+
r\left(u_{j-1}^n-2u_j^n+u_{j+1}^n\right),
\label{eq:explicit-heat-method}
\end{equation}

where

\begin{equation}
r=\frac{\alpha\Delta t}{(\Delta x)^2}.
\label{eq:stability-ratio}
\end{equation}

The value \(r\) is often called the stability ratio. For the explicit method, the time step must be chosen carefully. If \(r\) is too large, the numerical solution may become unstable.

\section{Algorithm}
\label{sec:computational-algorithm}

The method can be summarized as follows.

\begin{enumerate}
  \item Choose the number of spatial grid points and the final time.
  \item Compute \(\Delta x\), \(\Delta t\), and \(r\).
  \item Set the initial condition at each grid point.
  \item Apply the boundary conditions.
  \item Step forward in time using Equation~\ref{eq:explicit-heat-method}.
  \item Store or plot the approximate solution.
\end{enumerate}

When writing for accessibility, it is helpful to describe the algorithm in prose before showing code. The prose gives readers the purpose of the code and the main computational steps.

\section{Python Example}
\label{sec:python-example}

\noindent\textbf{Code Example 1.}
The following short Python script applies the explicit finite difference method to the one-dimensional heat equation. The code is included as real text rather than as a screenshot.

\begin{verbatim}
import numpy as np
import matplotlib.pyplot as plt

# Parameters
alpha = 1.0
x_min = 0.0
x_max = 1.0
final_time = 0.05

num_x = 51
dx = (x_max - x_min) / (num_x - 1)
dt = 0.4 * dx**2 / alpha
num_steps = int(final_time / dt)

r = alpha * dt / dx**2

# Grid and initial condition
x = np.linspace(x_min, x_max, num_x)
u = np.sin(np.pi * x)

# Boundary conditions
u[0] = 0.0
u[-1] = 0.0

# Time stepping
for n in range(num_steps):
    u_new = u.copy()

    for j in range(1, num_x - 1):
        u_new[j] = u[j] + r * (u[j - 1] - 2*u[j] + u[j + 1])

    u_new[0] = 0.0
    u_new[-1] = 0.0
    u = u_new

# Plot the final approximation
plt.plot(x, u)
plt.xlabel("x")
plt.ylabel("temperature")
plt.title("Approximate solution of the heat equation")
plt.show()
\end{verbatim}

The script first defines the physical and numerical parameters. It then creates the spatial grid, applies the initial condition, and advances the solution in time. The inner loop applies the finite difference formula from Equation~\ref{eq:explicit-heat-method} at each interior grid point.

\section{MATLAB Example}
\label{sec:matlab-example}

\noindent\textbf{Code Example 2.}
The following MATLAB script implements the same numerical method. Including both examples is useful in a template because some students work primarily in Python, while others work in MATLAB.

\begin{verbatim}
% Parameters
alpha = 1.0;
x_min = 0.0;
x_max = 1.0;
final_time = 0.05;

num_x = 51;
dx = (x_max - x_min) / (num_x - 1);
dt = 0.4 * dx^2 / alpha;
num_steps = floor(final_time / dt);

r = alpha * dt / dx^2;

% Grid and initial condition
x = linspace(x_min, x_max, num_x);
u = sin(pi * x);

% Boundary conditions
u(1) = 0.0;
u(end) = 0.0;

% Time stepping
for n = 1:num_steps
    u_new = u;

    for j = 2:num_x-1
        u_new(j) = u(j) + r * (u(j-1) - 2*u(j) + u(j+1));
    end

    u_new(1) = 0.0;
    u_new(end) = 0.0;
    u = u_new;
end

% Plot the final approximation
plot(x, u, 'LineWidth', 2)
xlabel('x')
ylabel('temperature')
title('Approximate solution of the heat equation')
grid on
\end{verbatim}

The MATLAB and Python versions implement the same mathematical method. The indexing differs because MATLAB arrays start at index \(1\), while Python arrays start at index \(0\).

\section{A Small Numerical Summary}
\label{sec:numerical-summary}

Table~\ref{tab:computational-summary} shows a small example of how computational results might be summarized. The table uses real text and column headers rather than a screenshot of program output.

\begin{table}[htbp]
\caption{Sample numerical summary for the explicit heat equation computation.}
\label{tab:computational-summary}
\centering
\begin{tabular}{llll}
\toprule
Grid points & Time step & Stability ratio \(r\) & Qualitative behavior \\
\midrule
\(25\)  & \(6.94\cdot 10^{-4}\) & \(0.40\) & Stable \\
\(51\)  & \(1.60\cdot 10^{-4}\) & \(0.40\) & Stable \\
\(101\) & \(4.00\cdot 10^{-5}\) & \(0.40\) & Stable \\
\bottomrule
\end{tabular}
\end{table}

A table like this is usually easier to read than a screenshot of a terminal window. If a table becomes too large, the student should consider using a `longtable` or moving detailed results to an appendix.

\section{A Simple Workflow Diagram}
\label{sec:workflow-diagram}

Figure~\ref{fig:computational-workflow} gives a schematic workflow for a computational experiment. This is the kind of figure that might appear in a computer science, engineering, or data science dissertation.

\begin{figure}[htbp]
\centering
\begin{tikzpicture}[
  node distance=1.1cm,
  >=Latex,
  box/.style={
    draw,
    rounded corners=3pt,
    minimum width=3.2cm,
    minimum height=0.8cm,
    align=center
  },
  arrow/.style={->, thick},
  alt={A vertical workflow diagram with five boxes. The boxes read define model, choose grid, run simulation, check stability, and analyze output. Arrows point downward from each step to the next.}
]
  \node[box] (model) {Define model};
  \node[box, below=of model] (grid) {Choose grid};
  \node[box, below=of grid] (run) {Run simulation};
  \node[box, below=of run] (check) {Check stability};
  \node[box, below=of check] (analyze) {Analyze output};

  \draw[arrow] (model) -- (grid);
  \draw[arrow] (grid) -- (run);
  \draw[arrow] (run) -- (check);
  \draw[arrow] (check) -- (analyze);
\end{tikzpicture}
\caption{A simple workflow for a computational experiment. The researcher defines the model, chooses a grid, runs the simulation, checks stability, and analyzes the output.}
\label{fig:computational-workflow}
\end{figure}

\section{Accessibility Notes for Code and Algorithms}
\label{sec:code-accessibility-notes}

Code should not be included as a screenshot. Screenshots of code are difficult to search, enlarge, copy, or interpret with assistive technology. Instead, use a real text block, and explain the purpose of the code before or after the block.

For accessible computational writing, authors should follow these practices.

\begin{enumerate}
  \item Introduce the mathematical or computational problem before showing code.
  \item Keep code examples short when possible.
  \item Explain important variables in prose.
  \item Avoid using color alone to communicate meaning in code or figures.
  \item Do not paste screenshots of terminal output, tables, or source code.
  \item Put long scripts in an appendix or a repository, and include only essential excerpts in the main text.
\end{enumerate}

These habits help readers understand both the computation and the reason it matters.
\chapter{CONCLUSIONS}
\label{chap:six}

This chapter summarizes the main findings of the dissertation, discusses the broader conclusions suggested by the work, and identifies possible directions for future study. The goal is to bring together the major themes developed in the earlier chapters and to explain how they contribute to the overall purpose of the research.

\section{Summary of Results}
\label{sec:summ}

The results of this dissertation address the central questions introduced at the beginning of the study. Through the methods and analysis presented in the previous chapters, the work identifies several important patterns, relationships, and outcomes. These results provide evidence that the main topic can be understood through a structured approach that connects the background material, the research design, and the interpretation of the findings.

Overall, the study shows that the problem under consideration has meaningful features that can be described, compared, and evaluated. The results also suggest that the framework used in this dissertation is useful for organizing the key ideas and for explaining how the individual parts of the study relate to the broader research question.

\section{Overarching Conclusions}
\label{sec:conc}

The conclusions drawn from this work emphasize the importance of careful analysis, clear organization, and thoughtful interpretation. The dissertation demonstrates that the topic is significant not only because of its immediate results, but also because of the larger questions it raises. By examining the problem from multiple perspectives, the study contributes to a more complete understanding of the subject.

A major conclusion is that the ideas developed throughout the dissertation are connected in ways that support the central argument of the work. The findings reinforce the value of the approach taken here and suggest that similar methods may be useful in related contexts. At the same time, the conclusions acknowledge that the study has limits and that additional work is needed to build on these results.

\section{Future Work}
\label{sec:future}

Future work may extend this dissertation in several directions. One possible direction is to apply the methods used here to a broader collection of examples or cases. This would make it possible to test the strength of the conclusions in additional settings and to determine whether the observed patterns remain consistent.

Another direction is to refine the framework introduced in this study. Further research could examine related questions in greater detail, consider alternative methods, or compare the present approach with other approaches. These extensions would help clarify the scope of the results and may lead to new insights about the topic.

Finally, future work may explore practical applications of the findings. By connecting the conclusions of this dissertation to new problems, future studies can continue to develop the ideas presented here and strengthen their usefulness for both research and practice.

%% -------------------------------------------------------------------------- %%
%% Bibliography
%% -------------------------------------------------------------------------- %%
%% Run BibTeX on the main file name:
%%
%%   lualatex dissertation-main-accessible
%%   bibtex dissertation-main-accessible
%%   lualatex dissertation-main-accessible
%%   lualatex dissertation-main-accessible

\bibliography{bibliography}
\bibliographystyle{apalike}

%% -------------------------------------------------------------------------- %%
%% Appendices
%% -------------------------------------------------------------------------- %%
%% \restoregeometry resets the page layout to the original thesis margins after
%% a temporary change made with \newgeometry{...}. It is included here as a
%% safeguard in case a chapter used temporary geometry settings.

\restoregeometry

%% \appendix tells LaTeX that the remaining chapters are appendices.
%% After this point, chapter numbering usually switches from numbers to letters:
%% Chapter 1, Chapter 2, ... become Appendix A, Appendix B, ...

\appendix

%% Include appendix files below, or delete/comment them out.

\include{Appendix-A/Appendix-A}
\include{Appendix-B/Appendix-B}

%% This second \restoregeometry is another safeguard. For best practice, any
%% temporary geometry change should be restored inside the same file where it was
%% changed, but this helps prevent accidental margin changes from carrying into
%% the back matter.

\restoregeometry

%% -------------------------------------------------------------------------- %%
%% Back matter
%% -------------------------------------------------------------------------- %%
%% \backmatter marks final material after the appendices. It is safe to keep,
%% even if the template does not currently include additional back matter pages.

\backmatter

\end{document}
